\section{方法对比与设计权衡}

\subsection{三种似然分数近似的对比}

\begin{table}[H]
\centering
\begin{tabular}{p{0.2\textwidth}p{0.25\textwidth}p{0.2\textwidth}p{0.2\textwidth}}
\toprule
\textbf{方法} & \textbf{$p(x_0|x_t)$ 的近似} & \textbf{$\mathcal{A}$ 的要求} & \textbf{近似精度} \\
\midrule
最粗暴近似 & $p(y|x_t) = p(y|x_0)$ & 任意 & 最差 \\
DPS & $\delta(x_0 - \hat{x}_0)$ & 可微即可 & 较差 \\
伪逆引导 & $\mathcal{N}(\hat{x}_0, r_t^2 I)$ & 必须线性 & 相对较好 \\
\bottomrule
\end{tabular}
\caption{三种似然分数近似方法的对比}
\end{table}

\begin{importantbox}{近似精度 vs. 适用范围的权衡}
这三种方法体现了一个核心设计权衡：
\begin{itemize}
    \item 更精确的近似（如伪逆引导）需要更强的假设（$\mathcal{A}$ 必须是线性的），但理论上更可靠。
    \item 更粗糙的近似（如 DPS）对 $\mathcal{A}$ 的要求更低（只需可微），适用范围更广，但近似误差更大。
\end{itemize}
没有一种方法在所有情况下都最优，选择取决于具体的应用场景和前向映射的性质。
\end{importantbox}

\subsection{所有方法的共同局限}

\begin{warningbox}{近似引导的固有局限}
无论采用哪种近似，这些推理时引导方法都有以下共同局限：
\begin{enumerate}
    \item \textbf{不保证收敛到正确的后验分布}：所有近似都引入了误差，生成结果可能不完全满足约束条件。
    \item \textbf{需要调节超参数}：引导强度（step size / guidance scale）需要手动调整，不同任务和不同数据可能需要不同的值。
    \item \textbf{计算开销}：DPS 需要对 Tweedie 估计进行反向传播，增加了每步的计算成本。
    \item \textbf{高噪声时间步近似更差}：在早期（高噪声）时间步，$\hat{x}_0$ 的估计质量很差，导致引导信号不可靠。
\end{enumerate}
\end{warningbox}

\subsection{本章小结}

不同的似然分数近似方法在精度和适用范围之间形成了一个谱系。实际应用中，需要根据前向映射的性质、计算预算和结果质量要求来选择合适的方法。


